\section{Conclusion} \label{sec:conslusion}
In this paper, we introduced \textsc{QuaRK}, an end-to-end quantum--classical learner for time-series prediction that couples a contractive quantum reservoir featurizer with a kernel-based RKHS readout. Our pipeline operates on windowed inputs and leverages a Johnson--Lindenstrauss projection to decouple the required quantum resources from the ambient data dimension. To enhance expressiveness while keeping the quantum component hardware-realistic, \textsc{QuaRK} relies on spatial multiplexing through multiple sub-reservoirs, and forms compact feature vectors by estimating families of $k$-local Pauli observables using classical-shadow tomography; these features are then learned with kernel ridge regression, where regularization provides an explicit and interpretable control of readout complexity.

On the learning-theoretic side, we connected these design choices to finite-sample performance under weak dependence. We established an effective interpolation regime, showing that once the embedding preserves enough geometry (via the prescribed projection/qubit scaling), relaxing the RKHS constraint yields a class rich enough to drive the empirical risk to (numerically) zero. We then provided a PAC-style generalization result for strided windows extracted from a stationary $\beta$-mixing process, recovering i.i.d.-like rates up to an explicit dependence penalty; empirically, our synthetic $\beta$-mixing VARMA benchmarks confirm the interpolation transition and the expected test-error decay when sweeping the number of training windows with a fixed featurizer and kernel.

Several directions remain to further consolidate \textsc{QuaRK} as a practical and hardware-relevant methodology. A natural next step is to sharpen the analysis under realistic noise and finite-shot effects, and to make the tradeoffs between measurement budget, observable locality $k$, multiplexing, and accuracy more explicit. Finally, more systematic tuning of the design choices (projection dimension/qubit count, sub-reservoir count, observable families, and Mat\'ern hyper-parameters) and validation on higher-dimensional, longer-memory, and nonstationary real-world time series will clarify when each component of the pipeline (projection, contraction, multiplexing, and kernel readout) yields the strongest gains.
